% ---
% titre: Calculs de ppmc et pgdc
% theme: NO
% chapitre: Nombres naturels et décimaux
% sous_chapitre: PPMC et PGDC
% niveau: [9e]
% type: EVAL
% tags: [c-ppmc, p-question-TS]
% difficulte: 1
% corrige: true
% ---

\question
Détermine:

\begin{parts}

\vspace{20pt}\part[3]
le ppmc de 14 et 40.

\vspace{20pt}{\footnotesize\raggedleft Espace pour ta démarche et tes calculs\par}
	\begin{solutionorgrid}[130pt]
		\begin{minipage}[t]{0.1\linewidth}
		\vspace{0pt}
		\begin{tabularx}{\linewidth}{>{\raggedleft\arraybackslash}X|X}
		14&2\\
		7&7\\
		1&\\
		\end{tabularx}
		\end{minipage}%
		\hfill
		\begin{minipage}[t]{0.1\linewidth}
		\vspace{0pt}
		\begin{tabularx}{\linewidth}{>{\raggedleft\arraybackslash}X|X}
		40&2\\
		20&2\\
		10&2\\
		5&5\\
		1&\\
		\end{tabularx}
		\end{minipage}%
		\hfill
		\begin{minipage}[t]{0.6\linewidth}
		\vspace{0pt}
		$14 = 2 \cdot 7$ \quad / \quad $40 = 2^3 \cdot 5$
		
		ppmc(14;40) $= 2^3 \cdot 5 \cdot 7 = 280$
		
		\vspace{20pt}
		1pt pour chaque décomposition ou liste de multiples
		
		1pt pour le ppmc
		\end{minipage}%
	\end{solutionorgrid}

\vspace{20pt}\part[3]
le pgdc de 144 et 270.

\vspace{20pt}{\footnotesize\raggedleft Espace pour ta démarche et tes calculs\par}
	\begin{solutionorgrid}[130pt]
		\begin{minipage}[t]{0.15\linewidth}
		\vspace{0pt}
		\begin{tabularx}{\linewidth}{>{\raggedleft\arraybackslash}X|X}
		144&2\\
		72&2\\
		36&2\\
		18&2\\
		9&3\\
		3&3\\
		1&\\
		\end{tabularx}
		\end{minipage}%
		\hfill
		\begin{minipage}[t]{0.15\linewidth}
		\vspace{0pt}
		\begin{tabularx}{\linewidth}{>{\raggedleft\arraybackslash}X|X}
		270&2\\
		135&3\\
		45&3\\
		15&3\\
		5&5\\
		1&\\
		\end{tabularx}
		\end{minipage}%
		\hfill
		\begin{minipage}[t]{0.6\linewidth}
		\vspace{0pt}
		$144 = 2^4 \cdot 3^2$ \quad / \quad $270 = 2 \cdot 3^3 \cdot 5$
		
		pgdc(144;270) $= 2 \cdot 3^2  = 18$
		
		\vspace{20pt}
		1pt pour chaque décomposition ou liste de diviseurs
		
		1pt pour le pgdc
		\end{minipage}
	\end{solutionorgrid}
	
\vspace{20pt}\part[4]
le ppmc de 45, 75 et 100.

\vspace{20pt}{\footnotesize\raggedleft Espace pour ta démarche et tes calculs\par}
	\begin{solutionorgrid}[150pt]
		\begin{minipage}[t]{0.13\linewidth}
		\vspace{0pt}
		\begin{tabularx}{\linewidth}{>{\raggedleft\arraybackslash}X|X}
		45&3\\
		15&3\\
		5&5\\
		1&\\
		\end{tabularx}
		\end{minipage}%
		\hfill
		\begin{minipage}[t]{0.13\linewidth}
		\vspace{0pt}
		\begin{tabularx}{\linewidth}{>{\raggedleft\arraybackslash}X|X}
		75&3\\
		25&5\\
		5&5\\
		1&\\
		\end{tabularx}
		\end{minipage}%
		\hfill
		\begin{minipage}[t]{0.13\linewidth}
		\vspace{0pt}
		\begin{tabularx}{\linewidth}{>{\raggedleft\arraybackslash}X|X}
		100&2\\
		50&2\\
		25&5\\
		5&5\\
		1&\\
		\end{tabularx}
		\end{minipage}%
		\hfill
		\begin{minipage}[t]{0.6\linewidth}
		\vspace{0pt}
		$45 = 3^2 \cdot 5$ \quad / \quad $75 = 3 \cdot 5^2$ \quad / \quad $100 = 2^2 \cdot 5^2$
		
		pgdc(45;75;100) $= 2^2 \cdot 3^2 \cdot 5^2  = 900$
		
		\vspace{20pt}
		1pt pour chaque décomposition ou liste de multiples
		
		1pt pour le ppmc
		\end{minipage}
	\end{solutionorgrid}
\end{parts}

